\documentclass[11pt]{article}

\usepackage[T1]{fontenc}
\usepackage[utf8]{inputenc}
\usepackage{lmodern}
\usepackage[margin=1in]{geometry}
\usepackage{amsmath,amssymb,amsthm,mathtools}
\usepackage{booktabs,array}
\usepackage{xcolor}
\usepackage{hyperref}
\usepackage{xurl}
\usepackage{microtype}
\newcommand{\OPHPaperReleaseID}{r2011}
\newcommand{\OPHPaperReleaseDate}{August 1, 2026}
\newcommand{\OPHPaperReleaseBanner}{%
\begin{center}
\small\textbf{Paper release:} \texttt{\OPHPaperReleaseID}\quad\textbf{Released:} \OPHPaperReleaseDate
\end{center}
}
\newcommand{\PaperReleaseID}{\OPHPaperReleaseID}
\newcommand{\PaperReleaseDate}{\OPHPaperReleaseDate}
\newcommand{\PaperReleaseBanner}{%
\begin{center}
\small\textbf{Paper release:} \texttt{\PaperReleaseID}\quad\textbf{Released:} \PaperReleaseDate
\end{center}
}


\hypersetup{
  colorlinks=true,
  linkcolor=blue,
  citecolor=blue,
  urlcolor=blue,
  pdftitle={The de Sitter Time-Advance Sign from a Finite Screen with Fixed Capacity},
  pdfauthor={Bernhard Mueller}
}
\setlength{\parindent}{0pt}
\setlength{\parskip}{0.55em}
\setlength{\emergencystretch}{3em}
\renewcommand{\arraystretch}{1.12}

\newtheorem{theorem}{Theorem}[section]
\newtheorem{proposition}[theorem]{Proposition}
\newtheorem{lemma}[theorem]{Lemma}
\newtheorem{assumption}[theorem]{Assumption}
\theoremstyle{remark}
\newtheorem{remark}[theorem]{Remark}

\newcommand{\R}{\mathbb R}
\newcommand{\one}{\mathbf 1}
\newcommand{\Hess}{\operatorname{Hess}}
\newcommand{\KL}{D_{\mathrm{KL}}}

\title{\textbf{The de Sitter Time-Advance Sign\\
from a Finite Screen with Fixed Capacity}}
\author{Bernhard Mueller\thanks{Corresponding author:
\href{mailto:bernhard@floatingpragma.ai}{\texttt{bernhard@floatingpragma.ai}}}\\
Pragma Research Inc., Washington, United States}
\date{\OPHPaperReleaseDate}

\begin{document}
\maketitle
\OPHPaperReleaseBanner

\begin{abstract}
Negative de Sitter shocks produce a time advance and an increasing
regularized out-of-time-ordered correlator. This sign conflicts with the
positivity and cyclicity required by a proposed static-patch trace. Under a
horizon-capacity dictionary, an exact fixed-budget law on a finite screen
supplies a conditional sign mechanism. Maximizing generalized entropy over sector probabilities gives
\(S_{\mathrm{gen}}^{\max}=\log M\). Uniform transfer of a fraction \(f\) of
the horizon capacity to an observer changes both this maximum and the
uniformly normalized area observable by \(\log(1-f)<0\). The unperturbed
horizon is therefore a one-sided boundary maximum along the transfer
direction. The exact gradient and Hessian show that this statement cannot be
replaced by an interior-maximum claim: at fixed horizon capacity the tangent
Hessian is positive, while the homogeneous curvature is negative. For an
icosahedral screen, a normalized finite shock spectrum follows conditionally
from exact gauge transport of the rotation triplet and a nearest-neighbour
Laplacian identification. The calculation does not derive the coefficient or
gravitational attachment. Its exact outputs are the pure-de-Sitter normalization, finite entropy
maximum, transfer law, analytic curvature, and line-graph spectrum identity.
The argument supplies no static-patch trace.
\end{abstract}

\noindent\textbf{Keywords:} de Sitter horizon, shock wave, static patch,
finite screen, generalized entropy, icosahedral graph

\section{The sign problem}

Chen, Stanford, Tang, and Yang study a static-patch proposal in which a
Euclidean gravitational path integral with an observer worldline computes a
Hilbert-space trace~\cite{chenStanfordTangYang2026}. Their shock-wave
calculation includes observer recoil and gravitational backreaction. The
resulting out-of-time-ordered correlator conflicts with cyclicity and
positivity: its leading regularized contribution increases. Geometrically,
the negative shock advances the signal.

The finite-screen calculation below isolates a capacity mechanism for this
sign. Three levels are kept separate throughout:
\begin{enumerate}
\item exact entropy and graph identities for a finite sector model;
\item an analytic relaxation of integer sector dimensions for stating
gradients and curvatures;
\item the physical dictionary identifying screen capacity transfer with a
gravitational horizon shock.
\end{enumerate}
The first two levels are mathematical statements. The third is a conditional
attachment. In particular, a finite capacity premise does not evade the trace
obstruction in Ref.~\cite{chenStanfordTangYang2026}. Finite-dimensionality is
part of the premise tested there.

\section{The smooth de Sitter normalization}

Write the horizon shock equation in spacetime dimension \(d\) as
\begin{equation}
\left(-\nabla^2_{S^{d-2}}-\mu^2\right)X^+
=8\pi G\,T^-,
\qquad
\mu^2=\frac{d-2}{2}|f'(r_c)|r_c=(d-2)\kappa r_c .
\label{eq:shock}
\end{equation}
For pure de Sitter space of radius \(L\),
\[
f(r)=1-\frac{r^2}{L^2},
\qquad
r_c=L,
\qquad
\kappa=\frac1L.
\]

\begin{lemma}[Radius-independent shock mass]
\label{lem:mu2}
For pure de Sitter space,
\[
\mu^2=d-2=\lambda_{\ell=1}\!\left(S^{d-2}\right).
\]
Consequently, \(\mu^2\) is independent of \(L\) and of the cosmological
constant.
\end{lemma}

\begin{proof}
The product \(\kappa r_c\) equals one. Equation~\eqref{eq:shock} therefore
gives \(\mu^2=d-2\). Scalar spherical harmonics on \(S^{d-2}\) have
\[
\lambda_\ell=\ell(\ell+d-3),
\]
so \(\lambda_1=d-2\).
\end{proof}

The smooth \(\ell=1\) shock is generated by a de Sitter isometry and is pure
gauge, with the corresponding source component removed by the gravitational
constraint~\cite{chenStanfordTangYang2026}. This explains the zero of the
operator in Eq.~\eqref{eq:shock}. Transferring that normalization to a
discrete screen requires a separate premise, stated in
Section~\ref{sec:spectrum}.

\section{Finite entropy and the capacity-transfer direction}

\subsection{The exact entropy maximum}

Let a finite screen have positive integer sector dimensions
\(d_1,\ldots,d_n\), and let \(p_i\) be a probability distribution over the
sectors. Set
\[
M=\sum_{i=1}^n d_i,
\qquad
S_{\mathrm{gen}}(p,d)
=-\sum_{i=1}^n p_i\log p_i+\sum_{i=1}^n p_i\log d_i.
\]

\begin{proposition}[Finite sector entropy maximum]
\label{prop:entropymax}
For fixed \(d_i\),
\[
p_i^\star=\frac{d_i}{M},
\qquad
\max_p S_{\mathrm{gen}}(p,d)=\log M.
\]
\end{proposition}

\begin{proof}
With \(q_i=d_i/M\),
\[
S_{\mathrm{gen}}(p,d)
=\log M-\sum_i p_i\log\frac{p_i}{q_i}
=\log M-\KL(p\Vert q).
\]
Nonnegativity of relative entropy proves the result.
\end{proof}

Thus a logarithmic capacity \(N=\log M_0\) agrees exactly with the extremal
generalized entropy when the closed screen carries total dimension \(M_0\).
This is a maximization over \(p\), with the sector dimensions held fixed. It
does not classify extrema with respect to the \(d_i\).

\subsection{Exact gradient and Hessian}

At the maximizing distribution, the expectation of the logarithmic sector
operator is
\begin{equation}
\mathcal A(d)
:=\sum_{i=1}^n\frac{d_i}{M}\log d_i .
\label{eq:areaobservable}
\end{equation}
The dimensions are integers in the finite model. Extend them to positive
real variables solely for the following differential calculation.

\begin{proposition}[Analytic curvature of the sector observable]
\label{prop:hessian}
The exact gradient of Eq.~\eqref{eq:areaobservable} is
\begin{equation}
\frac{\partial\mathcal A}{\partial d_i}
=\frac{\log d_i+1-\mathcal A}{M}.
\label{eq:gradient}
\end{equation}
Writing \(g_i=\partial_i\mathcal A\), its Hessian is
\begin{equation}
\frac{\partial^2\mathcal A}{\partial d_i\partial d_j}
=\frac{\delta_{ij}}{M d_i}-\frac{g_i+g_j}{M}.
\label{eq:hessiangeneral}
\end{equation}
At the symmetric point \(d_i=d\),
\begin{equation}
\Hess\mathcal A
=\frac1{n d^2}\left(I-\frac2nJ\right),
\label{eq:hessiansymmetric}
\end{equation}
where \(J=\one\one^T\). Its eigenvalue on \(\one\) is
\(-1/(nd^2)\), and every vector orthogonal to \(\one\) has eigenvalue
\(+1/(nd^2)\).
\end{proposition}

\begin{proof}
Differentiate \(\mathcal A=M^{-1}\sum_i d_i\log d_i\) to obtain
Eq.~\eqref{eq:gradient}. A second derivative gives
Eq.~\eqref{eq:hessiangeneral}. At \(d_i=d\), one has
\(M=nd\), \(\mathcal A=\log d\), and \(g_i=1/(nd)\), which yields
Eq.~\eqref{eq:hessiansymmetric}.
\end{proof}

Two consequences prevent a false extremum statement. First,
\(\nabla\mathcal A=\one/(nd)\) at the symmetric point, so the point is not
stationary in the full positive orthant. Second, a fixed-horizon-capacity
variation obeys \(\one^T\delta d=0\), and hence
\[
\delta d^T(\Hess\mathcal A)\delta d
=\frac{\|\delta d\|^2}{nd^2}>0
\]
for every nonzero tangent variation. The symmetric point is locally
minimizing along the fixed-\(M\) redistribution directions. The negative
eigenvalue belongs to the homogeneous direction, which changes \(M\).

\subsection{The one-sided budget maximum}

The de Sitter sign concerns transfer between the horizon and the observer,
rather than redistribution among horizon sectors. Let the closed total
capacity be \(M_0\). Assign \(c_{\mathrm{obs}}\geq0\) units in the finite
capacity ledger to the observer, leaving
\[
M_h=M_0-c_{\mathrm{obs}}
\]
on the horizon. Suppose the transfer depletes every horizon sector by the
same factor \(1-f\), where \(f=c_{\mathrm{obs}}/M_0\). For integer dimensions, restrict to
values for which the depleted dimensions remain integers. Then
\begin{equation}
S_{\mathrm{gen},h}^{\max}(f)-S_{\mathrm{gen},h}^{\max}(0)
=\log(1-f),
\label{eq:entropydrop}
\end{equation}
and
\begin{equation}
\mathcal A((1-f)d)-\mathcal A(d)
=\log(1-f).
\label{eq:areadrop}
\end{equation}
Both relations are exact.

\begin{theorem}[Finite transfer law and analytic one-sided maximum]
\label{thm:transfersign}
Every admissible positive integer transfer has a smaller horizon entropy
maximum and uniformly depleted area observable than the zero-transfer
allocation. On the positive-real interpolation
\(0\leq c_{\mathrm{obs}}<M_0\), the change is strictly decreasing and concave:
\[
\frac{d}{dc_{\mathrm{obs}}}\log(M_0-c_{\mathrm{obs}})
=-\frac1{M_0-c_{\mathrm{obs}}}<0,
\qquad
\frac{d^2}{dc_{\mathrm{obs}}^2}\log(M_0-c_{\mathrm{obs}})
=-\frac1{(M_0-c_{\mathrm{obs}})^2}<0.
\]
The allocation \(c_{\mathrm{obs}}=0\) is a one-sided boundary maximum on both
the admissible finite transfers and the analytic transfer path.
\end{theorem}

For transfers of \(2\%\), \(5\%\), and \(10\%\),
\[
\log(0.98)=-0.020203\ldots,\qquad
\log(0.95)=-0.051293\ldots,\qquad
\log(0.90)=-0.105361\ldots.
\]
The exact finite order and its analytic interpolation are the budget statements
relevant to the shock sign. Proposition~\ref{prop:hessian} supplies compatible
negative homogeneous curvature, while also showing why the transfer result
cannot be described as an interior maximum at fixed horizon \(M\).

\subsection{Conditional gravitational interpretation}

In the horizon coordinates used in
Ref.~\cite{chenStanfordTangYang2026}, the constant term of the shock operator
carries the sign of \(B'(0)/B(0)=2r'(0)/r_c\). For a cosmological horizon,
the transverse area decreases from the horizon toward the observer, and the
shock has the time-advance sign. Under the two physical identifications
\[
\text{closed screen capacity}=\text{fixed horizon--observer budget},
\qquad
\mathcal A=\frac{A_{\mathrm{horizon}}}{4G},
\]
Theorem~\ref{thm:transfersign} supplies the corresponding conditional sign
mechanism. Identifying the capacity coordinate \(c_{\mathrm{obs}}\) with the
observer mass used in the gravitational shock calculation is an additional
physical dictionary. The finite entropy theorem and transfer law establish
none of these identifications.

\section{Conditional icosahedral spectrum}
\label{sec:spectrum}

Consider the graph Laplacian of the twelve vertices of the icosahedron. Its
eigenspaces decompose under the alternating group \(A_5\) on five letters as
\[
\R^{12}\simeq\mathbf1\oplus\mathbf3\oplus\mathbf5\oplus\mathbf3',
\]
with raw Laplacian eigenvalues
\[
0,\qquad 5-\sqrt5,\qquad 6,\qquad 5+\sqrt5,
\]
respectively.

\begin{assumption}[Discrete gauge transport
\(\mathsf{DS\text{-}GAUGE}\)]
\label{ass:dsgauge}
The \(\mathbf3\) generated by the screen isometries consists of exact gauge
zero modes of the discrete shock operator, in direct analogy with the smooth
\(\ell=1\) isometry modes.
\end{assumption}

\begin{assumption}[Nearest-neighbour kinetic identification
\(\mathsf{DS\text{-}LAPLACIAN}\)]
\label{ass:dslaplacian}
The discrete shock kinetic operator is the nearest-neighbour combinatorial
Laplacian on the port graph, or its edge-sector line graph, with one positive
overall scale.
\end{assumption}

These two assumptions are independent. Assumption~\ref{ass:dsgauge} fixes the
scale by placing the \(\mathbf3\) at the smooth value \(2\).
Assumption~\ref{ass:dslaplacian} selects the finite operator whose remaining
eigenvalues are being compared.

\begin{proposition}[Conditional normalized shock spectrum]
\label{prop:spectrum}
Under Assumptions~\ref{ass:dsgauge} and \ref{ass:dslaplacian}, the normalized
port Laplacian and shock spectra are
\[
\begin{array}{c@{\qquad}c@{\qquad}c@{\qquad}c}
A_5\text{ block} & \text{multiplicity} &
\lambda_{\mathrm{norm}} & \lambda_{\mathrm{shock}}\\
\hline
\mathbf1 & 1 & 0 & -2\\
\mathbf3 & 3 & 2 & 0\\
\mathbf5 & 5 & 3+\dfrac3{\sqrt5} & 1+\dfrac3{\sqrt5}\\[3pt]
\mathbf3' & 3 & 3+\sqrt5 & 1+\sqrt5
\end{array}
\]
The invariant mode agrees exactly with the smooth value \(-2\). The
\(\mathbf5\) shock eigenvalue is \(1+3/\sqrt5=2.341640\ldots\), and the
additional \(\mathbf3'\) triplet lies at
\(1+\sqrt5=2\varphi=3.236067\ldots\).
\end{proposition}

\begin{proof}
Multiply the graph Laplacian by \(2/(5-\sqrt5)\), then subtract
\(\mu^2=2\). Direct simplification gives the displayed values.
\end{proof}

The value \(-2\) carries no enhancement relative to the continuum invariant
mode. Any different comparison obtained by leaving the graph scale unfixed
is a normalization artifact.

\subsection{Ports and edge sectors}

The logarithmic sector operator naturally suggests an edge carrier. The
low-lying spectrum is unchanged by this choice.

\begin{lemma}[Regular line-graph spectrum]
\label{lem:linegraph}
Let \(G\) be a \(k\)-regular graph with \(k\geq2\), \(n\) vertices, and \(m\)
edges. The
line graph \(L(G)\) has the Laplacian spectrum of \(G\), together with the
eigenvalue \(2k\) repeated \(m-n\) times.
\end{lemma}

\begin{proof}
Let \(B\) be the unoriented vertex-edge incidence matrix. Then
\[
BB^T=A_G+kI,
\qquad
B^TB=A_{L(G)}+2I.
\]
The nonzero spectra of \(BB^T\) and \(B^TB\) agree, and the spectrum of the
\(m\times m\) matrix \(B^TB\) contains \(m-n\) additional zeros relative to
the spectrum of the \(n\times n\) matrix \(BB^T\), counted with
multiplicity. Since \(L(G)\) is \((2k-2)\)-regular, converting adjacency
eigenvalues to Laplacian eigenvalues gives the claim.
\end{proof}

For the icosahedron, \(k=5\), \(n=12\), and \(m=30\). The edge-sector
spectrum is therefore
\[
\{0^1,(5-\sqrt5)^3,6^5,(5+\sqrt5)^3,10^{18}\}.
\]
It contains the full port spectrum plus a high block at \(10\). In
particular, the normalization-free ratio between the \(\mathbf5\) and
rotation-triplet Laplacian eigenvalues is
\[
\frac{\lambda_{\mathbf5}}{\lambda_{\mathbf3}}
=\frac6{5-\sqrt5}
=\frac{3(5+\sqrt5)}{10}
=2.170820\ldots
\]
on both carriers. The dodecahedral face graph gives
\(\varphi^2=2.618033\ldots\), while the smooth
\(\lambda_{\ell=2}/\lambda_{\ell=1}\) ratio is \(3\).

Lemma~\ref{lem:linegraph} closes the port-versus-edge ambiguity once the
nearest-neighbour Laplacian premise is accepted. It does not establish that
premise.

\section{The repair generator has a different domain}

Let \(X_r\) be the configuration space at fixed sector structure, with
measure \(\pi_r\), and let \(P_v\) be the single-site heat-bath conditional
expectation. The repair generator
\[
L^{\mathrm{rep}}=\sum_v(I-P_v)
\]
acts on
\[
K_r=L^2(X_r,\pi_r).
\]
Since each \(P_v\) is an orthogonal projection,
\[
\langle f,L^{\mathrm{rep}}f\rangle
=\sum_v\|(I-P_v)f\|^2.
\]
Consequently,
\[
\ker L^{\mathrm{rep}}=\bigcap_v\operatorname{Ran}P_v.
\]
Reducing this intersection to the constants requires a separate
irreducibility or intersection proof and is unnecessary for the scope
statement here.

A horizon shock changes the sector dimensions \(d_\alpha\) in
\[
L_C=\sum_\alpha(\log d_\alpha)P_\alpha .
\]
Such a change deforms the configuration space and lies outside
\(L^2(X_r,\pi_r)\) at fixed sector structure. The repair generator therefore
does not act on the sector-dimension shock. Its spectral gap neither excludes
nor supplies the time-advance mode. The capacity-transfer law in
Theorem~\ref{thm:transfersign} is a response input on the enlarged
sector-dimension space; no dynamical generator on that space is supplied.

\section{Claim boundary}

\begin{center}
\begin{tabular}{@{}p{0.39\textwidth}p{0.52\textwidth}@{}}
\toprule
Statement & Scope \\
\midrule
\(\mu^2=d-2=\lambda_{\ell=1}\) in pure de Sitter &
exact smooth-geometry identity \\
\(S_{\mathrm{gen}}^{\max}=\log M\) &
exact finite entropy maximum over sector probabilities \\
Uniform-depletion drop \(\log(1-f)\) &
exact finite law for admissible integer dimensions \\
Gradient and Hessian of \(\mathcal A(d)\) &
exact analytic-relaxation identity \\
One-sided horizon-budget maximum &
exact order on admissible finite transfers and exact analytic interpolation;
gravitational reading conditional \\
Port/edge low-spectrum equality &
exact regular line-graph theorem \\
Normalized icosahedral shock spectrum &
conditional on exact discrete gauge transport and the scaled
nearest-neighbour Laplacian identification \\
Shock coefficient and physical screen attachment &
not derived \\
Static-patch trace representation &
excluded from the claim; the cited trace obstruction stands \\
\bottomrule
\end{tabular}
\end{center}

The coefficient multiplying the shock response requires a generalized-entropy
normalization and a physical horizon-area dictionary on one common domain.
The finite calculation fixes neither. It also gives no Hilbert-space trace
whose correlators satisfy the conditions tested in
Ref.~\cite{chenStanfordTangYang2026}.

\subsection*{Closed tested classes and surviving alternatives}

The external trace obstruction concerns the proposal tested in
Ref.~\cite{chenStanfordTangYang2026}: a Euclidean static-patch path integral
with an observer worldline interpreted as a positive cyclic Hilbert-space
trace for the stated correlators. Horizon-screen descriptions without that
trace interpretation are outside the tested class.

The repair-domain statement concerns the declared fixed-sector
\(L^{\mathrm{rep}}\) on \(L^2(X_r,\pi_r)\). A separately constructed
dimension-changing generator is not excluded. The graph-spectrum statement
concerns the scaled nearest-neighbour icosahedral port Laplacian and its line
graph under exact gauge transport. Weighted, longer-range, refinement-limit,
and other source-derived kinetic operators are possible and require their
own spectra.

None of these boundaries disproves Observer-Patch Holography. They delimit
one trace interpretation, one fixed-sector generator, and one conditional
finite kinetic operator. The exact finite entropy, transfer, and line-graph
identities survive independently of those physical attachments.

\section{Conclusion}

A closed capacity budget has a distinguished transfer direction. Moving
capacity from a horizon screen to its observer changes the extremal entropy
and the uniformly normalized logarithmic area observable by
\(\log(1-f)\). The sign is negative for every allowed positive finite
transfer and throughout the positive-real interpolation, so the zero-observer
allocation is a one-sided boundary maximum. The analytic Hessian supports a
negative homogeneous curvature and a positive fixed-horizon-capacity tangent
curvature. Keeping both facts visible avoids an incorrect interior-maximum
argument.

The pure de Sitter normalization \(\mu^2=\lambda_{\ell=1}\) is exact. An
icosahedral spectrum follows only after exact discrete gauge transport and a
nearest-neighbour kinetic operator are assumed. The regular line-graph
identity then makes its low spectrum independent of choosing ports or edge
sectors. These results supply a conditional finite-screen mechanism for the
time-advance sign, with the ledger-to-mass map, coefficient, gravitational
attachment, and sector-dimension dynamics as separate premises.
The exact identities remain valid if those attachments fail; such a failure
would reject the corresponding physical interpretation inside its stated
class.

\begin{thebibliography}{9}

\bibitem{chenStanfordTangYang2026}
Y.~Chen, D.~Stanford, H.~Tang, and Z.~Yang,
\emph{Negative shocks versus static patch holography},
arXiv:2607.14042 [hep-th], 2026.
\url{https://arxiv.org/abs/2607.14042}

\end{thebibliography}

\end{document}
